\chapter{Arquitectura del Cuaderno de Examenes}

\begin{planningbox}{Objetivo}
El cuaderno de examenes se proyecta como un complemento de entrenamiento exigente. Debe servir
para dos usos distintos:
\begin{itemize}
  \item mini-examenes por seccion, de \(25\) a \(35\) minutos, para cerrar un apartado;
  \item examenes por bloques, de \(90\) a \(120\) minutos, para evaluar transferencia,
  eleccion de metodo y resistencia cognitiva.
\end{itemize}
\end{planningbox}

\begin{exambrief}{Estructura base}
Cada mini-examen por seccion debe incluir:
\begin{enumerate}
  \item una tarea tecnica representativa del apartado;
  \item una tarea de interpretacion o modelizacion;
  \item una solucion completa con baremo.
\end{enumerate}

Cada examen por bloques debe incluir:
\begin{enumerate}
  \item al menos tres tecnicas diferentes;
  \item un problema con datos, tabla o grafico;
  \item un item de justificacion, no solo de calculo;
  \item una solucion docente cerrada.
\end{enumerate}
\end{exambrief}

\section{Inventario de mini-examenes por seccion}

\setlength{\LTpre}{0pt}
\setlength{\LTpost}{0pt}
\small
\begin{longtable}{@{}>{\raggedright\arraybackslash}p{0.10\textwidth}>{\raggedright\arraybackslash}p{0.22\textwidth}>{\raggedright\arraybackslash}p{0.25\textwidth}>{\raggedright\arraybackslash}p{0.18\textwidth}>{\raggedright\arraybackslash}p{0.09\textwidth}@{}}
\toprule
\textbf{ID} & \textbf{Seccion} & \textbf{Mini-examen propuesto} & \textbf{Bloque} & \textbf{Tiempo}\\
\midrule
\endfirsthead
\toprule
\textbf{ID} & \textbf{Seccion} & \textbf{Mini-examen propuesto} & \textbf{Bloque} & \textbf{Tiempo}\\
\midrule
\endhead
\multicolumn{5}{l}{\textbf{C01 --- Numeros reales y herramientas numericas}}\\
C01.S01 & Clasificacion de numeros y conjuntos numericos & Mini-examen de modelizacion algebraica y cierre exacto & B1 Numeros reales y algebra inicial & 25--35 min\\
C01.S02 & Decimales y fraccion generatriz & Mini-examen de modelizacion algebraica y cierre exacto & B1 Numeros reales y algebra inicial & 25--35 min\\
C01.S03 & Intervalos, union e interseccion & Mini-examen con representacion y lectura de solucion & B1 Numeros reales y algebra inicial & 25--35 min\\
C01.S04 & Valor absoluto como distancia & Mini-examen de modelizacion algebraica y cierre exacto & B1 Numeros reales y algebra inicial & 25--35 min\\
C01.S05 & Potencias y radicales & Mini-examen de modelizacion algebraica y cierre exacto & B1 Numeros reales y algebra inicial & 25--35 min\\
C01.S06 & Racionalizacion & Mini-examen de modelizacion algebraica y cierre exacto & B1 Numeros reales y algebra inicial & 25--35 min\\
C01.S07 & Logaritmos y propiedades operativas & Mini-examen de modelizacion algebraica y cierre exacto & B1 Numeros reales y algebra inicial & 25--35 min\\
\multicolumn{5}{l}{\textbf{C02 --- Polinomios y fracciones algebraicas}}\\
C02.S01 & Valor numerico y lectura de polinomios & Mini-examen de modelizacion algebraica y cierre exacto & B1 Numeros reales y algebra inicial & 25--35 min\\
C02.S02 & Productos notables y operaciones con polinomios & Mini-examen de modelizacion algebraica y cierre exacto & B1 Numeros reales y algebra inicial & 25--35 min\\
C02.S03 & Division de polinomios & Mini-examen de modelizacion algebraica y cierre exacto & B1 Numeros reales y algebra inicial & 25--35 min\\
C02.S04 & Teorema del resto, del factor y parametros & Mini-examen con decision algebraica y condicion adicional & B1 Numeros reales y algebra inicial & 25--35 min\\
C02.S05 & Raices y factorizacion & Mini-examen de modelizacion algebraica y cierre exacto & B1 Numeros reales y algebra inicial & 25--35 min\\
C02.S06 & Fracciones algebraicas: dominio y simplificacion & Mini-examen de modelizacion algebraica y cierre exacto & B1 Numeros reales y algebra inicial & 25--35 min\\
C02.S07 & Suma y resta de fracciones algebraicas & Mini-examen de modelizacion algebraica y cierre exacto & B1 Numeros reales y algebra inicial & 25--35 min\\
C02.S08 & Producto, cociente y operaciones combinadas & Mini-examen de modelizacion algebraica y cierre exacto & B1 Numeros reales y algebra inicial & 25--35 min\\
C02.S09 & Ecuaciones con fracciones algebraicas & Mini-examen de modelizacion algebraica y cierre exacto & B1 Numeros reales y algebra inicial & 25--35 min\\
\multicolumn{5}{l}{\textbf{C03 --- Ecuaciones y sistemas}}\\
C03.S01 & Ecuaciones racionales & Mini-examen de tecnica y transferencia & B2 Ecuaciones y sistemas & 25--35 min\\
C03.S02 & Ecuaciones irracionales & Mini-examen de tecnica y transferencia & B2 Ecuaciones y sistemas & 25--35 min\\
C03.S03 & Ecuaciones polinomicas reducibles & Mini-examen de tecnica y transferencia & B2 Ecuaciones y sistemas & 25--35 min\\
C03.S04 & Ecuaciones exponenciales & Mini-examen de tecnica y transferencia & B2 Ecuaciones y sistemas & 25--35 min\\
C03.S05 & Ecuaciones logaritmicas & Mini-examen de tecnica y transferencia & B2 Ecuaciones y sistemas & 25--35 min\\
C03.S06 & Eleccion de metodo en ecuaciones mixtas & Mini-examen de tecnica y transferencia & B2 Ecuaciones y sistemas & 25--35 min\\
C03.S07 & Sistemas lineales 2x2 & Mini-examen de traduccion, resolucion y comprobacion & B2 Ecuaciones y sistemas & 25--35 min\\
C03.S08 & Sistemas lineales 3x3 y clasificacion & Mini-examen de traduccion, resolucion y comprobacion & B2 Ecuaciones y sistemas & 25--35 min\\
C03.S09 & Problemas de planteamiento con sistemas & Mini-examen de traduccion, resolucion y comprobacion & B2 Ecuaciones y sistemas & 25--35 min\\
C03.S10 & Sistemas no lineales & Mini-examen de traduccion, resolucion y comprobacion & B2 Ecuaciones y sistemas & 25--35 min\\
C03.S11 & Sistemas exponenciales y logaritmicos & Mini-examen de traduccion, resolucion y comprobacion & B2 Ecuaciones y sistemas & 25--35 min\\
\multicolumn{5}{l}{\textbf{C04 --- Inecuaciones}}\\
C04.S01 & Inecuaciones lineales & Mini-examen con representacion y lectura de solucion & B2 Ecuaciones y sistemas & 25--35 min\\
C04.S02 & Inecuaciones de segundo grado & Mini-examen con representacion y lectura de solucion & B2 Ecuaciones y sistemas & 25--35 min\\
C04.S03 & Inecuaciones polinomicas & Mini-examen con representacion y lectura de solucion & B2 Ecuaciones y sistemas & 25--35 min\\
C04.S04 & Inecuaciones racionales & Mini-examen con representacion y lectura de solucion & B2 Ecuaciones y sistemas & 25--35 min\\
C04.S05 & Sistemas de inecuaciones en una variable & Mini-examen de traduccion, resolucion y comprobacion & B2 Ecuaciones y sistemas & 25--35 min\\
C04.S06 & Inecuaciones de dos variables y semiplanos & Mini-examen con representacion y lectura de solucion & B2 Ecuaciones y sistemas & 25--35 min\\
C04.S07 & Sistemas de inecuaciones en el plano & Mini-examen de traduccion, resolucion y comprobacion & B2 Ecuaciones y sistemas & 25--35 min\\
\multicolumn{5}{l}{\textbf{C05 --- Trigonometria}}\\
C05.S01 & Grados y radianes & Mini-examen con croquis obligatorio y justificacion geometrica & B3 Trigonometria y vectores & 25--35 min\\
C05.S02 & Razones trigonometricas en triangulos y poligonos & Mini-examen con croquis obligatorio y justificacion geometrica & B3 Trigonometria y vectores & 25--35 min\\
C05.S03 & Angulos notables y reduccion al primer cuadrante & Mini-examen con croquis obligatorio y justificacion geometrica & B3 Trigonometria y vectores & 25--35 min\\
C05.S04 & Calculo de razones a partir de una razon conocida & Mini-examen con croquis obligatorio y justificacion geometrica & B3 Trigonometria y vectores & 25--35 min\\
C05.S05 & Identidades y transformaciones trigonometricas & Mini-examen con croquis obligatorio y justificacion geometrica & B3 Trigonometria y vectores & 25--35 min\\
C05.S06 & Ecuaciones trigonometricas & Mini-examen con croquis obligatorio y justificacion geometrica & B3 Trigonometria y vectores & 25--35 min\\
C05.S07 & Resolucion de triangulos & Mini-examen con croquis obligatorio y justificacion geometrica & B3 Trigonometria y vectores & 25--35 min\\
C05.S08 & Problemas contextualizados de trigonometria & Mini-examen con croquis obligatorio y justificacion geometrica & B3 Trigonometria y vectores & 25--35 min\\
\multicolumn{5}{l}{\textbf{C06 --- Vectores}}\\
C06.S01 & Bases y dependencia en R2 & Mini-examen con croquis obligatorio y justificacion geometrica & B3 Trigonometria y vectores & 25--35 min\\
C06.S02 & Operaciones con vectores & Mini-examen con croquis obligatorio y justificacion geometrica & B3 Trigonometria y vectores & 25--35 min\\
C06.S03 & Modulo, argumento y vectores unitarios & Mini-examen con croquis obligatorio y justificacion geometrica & B3 Trigonometria y vectores & 25--35 min\\
C06.S04 & Producto escalar y angulo entre vectores & Mini-examen con croquis obligatorio y justificacion geometrica & B3 Trigonometria y vectores & 25--35 min\\
C06.S05 & Paralelismo, perpendicularidad y alineacion & Mini-examen con croquis obligatorio y justificacion geometrica & B3 Trigonometria y vectores & 25--35 min\\
C06.S06 & Construcciones con vectores y paralelogramos & Mini-examen con croquis obligatorio y justificacion geometrica & B3 Trigonometria y vectores & 25--35 min\\
\multicolumn{5}{l}{\textbf{C07 --- Geometria analitica y lugares geometricos}}\\
C07.S01 & Ecuaciones de la recta y vectores asociados & Mini-examen con croquis obligatorio y justificacion geometrica & B4 Geometria analitica & 25--35 min\\
C07.S02 & Problemas con parametros en rectas & Mini-examen con decision algebraica y condicion adicional & B4 Geometria analitica & 25--35 min\\
C07.S03 & Simetrias respecto de punto y recta & Mini-examen con croquis obligatorio y justificacion geometrica & B4 Geometria analitica & 25--35 min\\
C07.S04 & Posicion relativa de rectas & Mini-examen con croquis obligatorio y justificacion geometrica & B4 Geometria analitica & 25--35 min\\
C07.S05 & Rectas especiales por un punto & Mini-examen con croquis obligatorio y justificacion geometrica & B4 Geometria analitica & 25--35 min\\
C07.S06 & Distancias y angulos entre rectas & Mini-examen con croquis obligatorio y justificacion geometrica & B4 Geometria analitica & 25--35 min\\
C07.S07 & Poligonos y areas con coordenadas & Mini-examen con croquis obligatorio y justificacion geometrica & B4 Geometria analitica & 25--35 min\\
C07.S08 & Medianas, alturas y mediatrices & Mini-examen con croquis obligatorio y justificacion geometrica & B4 Geometria analitica & 25--35 min\\
C07.S09 & Bisectrices, ortocentro y lugares geometricos & Mini-examen con croquis obligatorio y justificacion geometrica & B4 Geometria analitica & 25--35 min\\
C07.S10 & Condiciones geometricas combinadas & Mini-examen con croquis obligatorio y justificacion geometrica & B4 Geometria analitica & 25--35 min\\
\multicolumn{5}{l}{\textbf{C08 --- Limites, continuidad y asintotas}}\\
C08.S01 & Dominio de funciones & Mini-examen de lectura analitica, grafica y conclusion & B5 Limites y continuidad & 25--35 min\\
C08.S02 & Lectura de graficas & Mini-examen de lectura analitica, grafica y conclusion & B5 Limites y continuidad & 25--35 min\\
C08.S03 & Limites basicos y laterales & Mini-examen de lectura analitica, grafica y conclusion & B5 Limites y continuidad & 25--35 min\\
C08.S04 & Limites en el infinito & Mini-examen de lectura analitica, grafica y conclusion & B5 Limites y continuidad & 25--35 min\\
C08.S05 & Indeterminaciones por factorizacion y racionalizacion & Mini-examen de lectura analitica, grafica y conclusion & B5 Limites y continuidad & 25--35 min\\
C08.S06 & Asintotas & Mini-examen de lectura analitica, grafica y conclusion & B5 Limites y continuidad & 25--35 min\\
C08.S07 & Continuidad y discontinuidades & Mini-examen de lectura analitica, grafica y conclusion & B5 Limites y continuidad & 25--35 min\\
C08.S08 & Funciones a trozos y representacion & Mini-examen de lectura analitica, grafica y conclusion & B5 Limites y continuidad & 25--35 min\\
C08.S09 & Continuidad con parametros & Mini-examen con decision algebraica y condicion adicional & B5 Limites y continuidad & 25--35 min\\
C08.S10 & Modelos de evolucion a largo plazo & Mini-examen de lectura analitica, grafica y conclusion & B5 Limites y continuidad & 25--35 min\\
C08.S11 & Ampliacion y estilo EBAU en limites y continuidad & Mini-examen de lectura analitica, grafica y conclusion & B5 Limites y continuidad & 25--35 min\\
\multicolumn{5}{l}{\textbf{C09 --- Derivadas y aplicaciones}}\\
C09.S01 & Reglas de derivacion & Mini-examen de lectura analitica, grafica y conclusion & B6 Derivadas y aplicaciones & 25--35 min\\
C09.S02 & Monotonia, extremos, curvatura e inflexion & Mini-examen de lectura analitica, grafica y conclusion & B6 Derivadas y aplicaciones & 25--35 min\\
C09.S03 & Recta tangente y recta normal en un punto & Mini-examen de lectura analitica, grafica y conclusion & B6 Derivadas y aplicaciones & 25--35 min\\
C09.S04 & Tangente con pendiente dada o paralela a una recta & Mini-examen de lectura analitica, grafica y conclusion & B6 Derivadas y aplicaciones & 25--35 min\\
C09.S05 & Problemas con parametros en derivadas & Mini-examen con decision algebraica y condicion adicional & B6 Derivadas y aplicaciones & 25--35 min\\
C09.S06 & Derivabilidad de funciones a trozos & Mini-examen de lectura analitica, grafica y conclusion & B6 Derivadas y aplicaciones & 25--35 min\\
C09.S07 & Optimizacion & Mini-examen de lectura analitica, grafica y conclusion & B6 Derivadas y aplicaciones & 25--35 min\\
C09.S08 & Ampliacion y selectividad en derivadas & Mini-examen de lectura analitica, grafica y conclusion & B6 Derivadas y aplicaciones & 25--35 min\\
\multicolumn{5}{l}{\textbf{C10 --- Estadistica, probabilidad e inferencia}}\\
C10.S01 & Datos bidimensionales y tablas de contingencia & Mini-examen con datos, interpretacion y decision bajo incertidumbre & B7 Estadistica, probabilidad e inferencia & 25--35 min\\
C10.S02 & Nubes de puntos, correlacion y causalidad & Mini-examen con datos, interpretacion y decision bajo incertidumbre & B7 Estadistica, probabilidad e inferencia & 25--35 min\\
C10.S03 & Regresion lineal y cuadratica & Mini-examen con datos, interpretacion y decision bajo incertidumbre & B7 Estadistica, probabilidad e inferencia & 25--35 min\\
C10.S04 & Muestreo, inferencia y tecnologia & Mini-examen con datos, interpretacion y decision bajo incertidumbre & B7 Estadistica, probabilidad e inferencia & 25--35 min\\
C10.S05 & Experimentos aleatorios y algebra de sucesos & Mini-examen con datos, interpretacion y decision bajo incertidumbre & B7 Estadistica, probabilidad e inferencia & 25--35 min\\
C10.S06 & Regla de Laplace y tecnicas de recuento & Mini-examen con datos, interpretacion y decision bajo incertidumbre & B7 Estadistica, probabilidad e inferencia & 25--35 min\\
C10.S07 & Probabilidad condicionada, total y Bayes & Mini-examen con datos, interpretacion y decision bajo incertidumbre & B7 Estadistica, probabilidad e inferencia & 25--35 min\\
\bottomrule
\end{longtable}


\section{Bloques de examen previstos}

\begin{databox}{Mapa de bloques}
\begin{center}
\begin{tabular}{ll}
\toprule
Bloque & Contenido \\
\midrule
\(B1\) & Numeros reales y algebra inicial (`C01-C02`) \\
\(B2\) & Ecuaciones, sistemas e inecuaciones (`C03-C04`) \\
\(B3\) & Trigonometria y vectores (`C05-C06`) \\
\(B4\) & Geometria analitica (`C07`) \\
\(B5\) & Limites, continuidad y asintotas (`C08`) \\
\(B6\) & Derivadas y aplicaciones (`C09`) \\
\(B7\) & Estadistica, probabilidad e inferencia (`C10`) \\
\bottomrule
\end{tabular}
\end{center}
\end{databox}

\chapter{Mini-Examenes Modelo por Seccion}

\section{Modelo C01.S07 --- Logaritmos y propiedades operativas}

\begin{exambrief}{Mini-examen de cierre}
Duracion sugerida: \(30\) minutos.
\begin{enumerate}
  \item Simplifica \(\log 50+\log 2-\log 5\).
  \item En una escala sonora se cumple \(L=10\log(I/I_0)\). Si un foco sonoro aumenta su
  intensidad cien veces, calcula el incremento de decibelios.
\end{enumerate}
\end{exambrief}

\begin{rubricbox}{C01.S07}
\begin{itemize}
  \item Item 1: \(4\) puntos.
  \item Item 2: \(6\) puntos.
\end{itemize}
\end{rubricbox}

\begin{shortanswer}{Mini-examen C01.S07}
\[
\log 50+\log 2-\log 5=\log 20,
\qquad
\Delta L=20\text{ dB}.
\]
\end{shortanswer}

\begin{fullsolution}{Mini-examen C01.S07}
\[
\log 50+\log 2-\log 5=\log\left(\frac{50\cdot 2}{5}\right)=\log 20.
\]
Si la intensidad se multiplica por \(100\),
\[
\Delta L=10\log(100)=10\cdot 2=20\text{ dB}.
\]
\end{fullsolution}

\section{Modelo C02.S04 --- Teorema del resto y parametros}

\begin{exambrief}{Mini-examen de cierre}
Duracion sugerida: \(30\) minutos.
\begin{enumerate}
  \item Halla \(a\) para que el polinomio \(P(x)=x^3-3x^2+ax+6\) sea divisible por \(x-2\).
  \item Con el valor obtenido, factoriza \(P(x)\) sabiendo que \(x-2\) es factor.
\end{enumerate}
\end{exambrief}

\begin{shortanswer}{Mini-examen C02.S04}
Se obtiene \(a=-1\) y
\[
P(x)=(x-2)(x^2-x-3).
\]
\end{shortanswer}

\begin{fullsolution}{Mini-examen C02.S04}
Para que \(x-2\) sea factor debe cumplirse \(P(2)=0\):
\[
8-12+2a+6=0
\Longrightarrow
2+2a=0
\Longrightarrow
a=-1.
\]
Entonces
\[
P(x)=x^3-3x^2-x+6.
\]
Dividiendo entre \(x-2\),
\[
\begin{array}{r|rrrr}
2 & 1 & -3 & -1 & 6\\
  &   & 2 & -2 & -6\\
\hline
  & 1 & -1 & -3 & 0
\end{array}
\]
y por tanto
\[
P(x)=(x-2)(x^2-x-3).
\]
\end{fullsolution}

\section{Modelo C03.S09 --- Problema de planteamiento con sistemas}

\begin{exambrief}{Mini-examen de cierre}
Duracion sugerida: \(35\) minutos.
\begin{enumerate}
  \item En una actividad se venden entradas general a \(14\) euros y reducida a \(9\) euros.
  Se venden \(260\) entradas y se recaudan \(2980\) euros. Calcula cuantas de cada tipo se
  vendieron.
  \item Comprueba la coherencia de la solucion en el contexto.
\end{enumerate}
\end{exambrief}

\begin{shortanswer}{Mini-examen C03.S09}
Se vendieron \(128\) entradas generales y \(132\) reducidas.
\end{shortanswer}

\begin{fullsolution}{Mini-examen C03.S09}
Si \(g\) es el numero de entradas generales y \(r\) el de reducidas:
\[
\begin{cases}
g+r=260,\\
14g+9r=2980.
\end{cases}
\]
Restando \(9(g+r)=2340\) a la segunda ecuacion:
\[
5g=640 \Longrightarrow g=128.
\]
Entonces
\[
r=260-128=132.
\]
La comprobacion es inmediata:
\[
128+132=260,\qquad 14\cdot 128+9\cdot 132=2980.
\]
\end{fullsolution}

\section{Modelo C04.S07 --- Sistemas de inecuaciones en el plano}

\begin{exambrief}{Mini-examen de cierre}
Duracion sugerida: \(35\) minutos.
\begin{enumerate}
  \item Representa la region definida por
  \[
  \begin{cases}
  x+y\leq 8,\\
  x\geq 1,\\
  y\geq 2.
  \end{cases}
  \]
  \item Halla sus vertices.
  \item Maximiza \(F(x,y)=3x+2y\) en dicha region.
\end{enumerate}
\end{exambrief}

\begin{shortanswer}{Mini-examen C04.S07}
Los vertices son \((1,2)\), \((1,7)\) y \((6,2)\). El maximo de \(F\) es \(22\) y se alcanza
en \((6,2)\).
\end{shortanswer}

\begin{fullsolution}{Mini-examen C04.S07}
La region queda limitada por la recta \(x+y=8\), la vertical \(x=1\) y la horizontal \(y=2\).
Sus vertices son:
\[
(1,2),\qquad (1,7),\qquad (6,2).
\]
Evaluamos:
\[
F(1,2)=7,\qquad F(1,7)=17,\qquad F(6,2)=22.
\]
Correccion: el valor mayor es \(22\), luego el maximo se alcanza en \((6,2)\).
\end{fullsolution}

\section{Modelo C05.S06 --- Ecuaciones trigonometricas}

\begin{exambrief}{Mini-examen de cierre}
Duracion sugerida: \(30\) minutos.
\begin{enumerate}
  \item Resuelve \(\sen x=\frac{\sqrt{3}}{2}\) en \([0,2\pi)\).
  \item Escribe la solucion general.
\end{enumerate}
\end{exambrief}

\begin{shortanswer}{Mini-examen C05.S06}
En \([0,2\pi)\),
\[
x=\frac{\pi}{3}\quad \text{o}\quad x=\frac{2\pi}{3}.
\]
La solucion general es
\[
x=\frac{\pi}{3}+2k\pi \quad \text{o}\quad x=\frac{2\pi}{3}+2k\pi,\qquad k\in\mathbb{Z}.
\]
\end{shortanswer}

\begin{fullsolution}{Mini-examen C05.S06}
El seno vale \(\frac{\sqrt{3}}{2}\) para el angulo de referencia \(\frac{\pi}{3}\). Como el
seno es positivo en los cuadrantes I y II, en \([0,2\pi)\) se obtiene
\[
x=\frac{\pi}{3},\qquad x=\pi-\frac{\pi}{3}=\frac{2\pi}{3}.
\]
Anadiendo vueltas completas:
\[
x=\frac{\pi}{3}+2k\pi \quad \text{o}\quad x=\frac{2\pi}{3}+2k\pi,\qquad k\in\mathbb{Z}.
\]
\end{fullsolution}

\section{Modelo C06.S04 --- Producto escalar y angulo}

\begin{exambrief}{Mini-examen de cierre}
Duracion sugerida: \(30\) minutos.
\begin{enumerate}
  \item Calcula el producto escalar de \(\vec u=(4,1)\) y \(\vec v=(2,5)\).
  \item Halla el angulo entre ambos vectores.
\end{enumerate}
\end{exambrief}

\begin{shortanswer}{Mini-examen C06.S04}
\[
\vec u\cdot \vec v=13,
\qquad
\cos\theta=\frac{13}{\sqrt{17}\sqrt{29}},
\]
de donde \(\theta\approx 54.2^\circ\).
\end{shortanswer}

\begin{fullsolution}{Mini-examen C06.S04}
\[
\vec u\cdot \vec v=4\cdot 2+1\cdot 5=13.
\]
Los modulos son
\[
\|\vec u\|=\sqrt{17},\qquad \|\vec v\|=\sqrt{29}.
\]
Por la formula del producto escalar,
\[
\cos\theta=\frac{\vec u\cdot \vec v}{\|\vec u\|\,\|\vec v\|}=
\frac{13}{\sqrt{17}\sqrt{29}}.
\]
Aplicando arco coseno:
\[
\theta\approx 54.2^\circ.
\]
\end{fullsolution}

\section{Modelo C07.S06 --- Distancia entre rectas}

\begin{exambrief}{Mini-examen de cierre}
Duracion sugerida: \(35\) minutos.
\begin{enumerate}
  \item Comprueba que las rectas \(r:2x-y+1=0\) y \(s:2x-y-5=0\) son paralelas.
  \item Calcula la distancia entre ellas.
\end{enumerate}
\end{exambrief}

\begin{shortanswer}{Mini-examen C07.S06}
Son paralelas porque comparten vector normal \((2,-1)\). La distancia es
\[
\frac{|1-(-5)|}{\sqrt{2^2+(-1)^2}}=\frac{6}{\sqrt{5}}.
\]
\end{shortanswer}

\begin{fullsolution}{Mini-examen C07.S06}
Ambas rectas tienen la forma
\[
2x-y+c=0,
\]
de modo que comparten el mismo vector normal \((2,-1)\) y son paralelas.
La distancia entre rectas paralelas \(ax+by+c_1=0\) y \(ax+by+c_2=0\) es
\[
d=\frac{|c_1-c_2|}{\sqrt{a^2+b^2}}.
\]
Aplicando la formula:
\[
d=\frac{|1-(-5)|}{\sqrt{2^2+(-1)^2}}=\frac{6}{\sqrt{5}}.
\]
\end{fullsolution}

\section{Modelo C08.S09 --- Continuidad con parametros}

\begin{exambrief}{Mini-examen de cierre}
Duracion sugerida: \(30\) minutos.
\begin{enumerate}
  \item Determina \(a\) para que la funcion
  \[
  f(x)=
  \begin{cases}
    ax+1, & x<2,\\
    x^2-1, & x\geq 2
  \end{cases}
  \]
  sea continua en \(x=2\).
  \item Indica el valor comun en el punto de empalme.
\end{enumerate}
\end{exambrief}

\begin{shortanswer}{Mini-examen C08.S09}
La continuidad exige \(2a+1=3\), luego \(a=1\). El valor comun es \(3\).
\end{shortanswer}

\begin{fullsolution}{Mini-examen C08.S09}
Para que haya continuidad en \(x=2\), el limite por la izquierda debe coincidir con el valor de
la rama derecha:
\[
\lim_{x\to 2^-}(ax+1)=2a+1,
\qquad
f(2)=2^2-1=3.
\]
Imponemos
\[
2a+1=3 \Longrightarrow 2a=2 \Longrightarrow a=1.
\]
Con ese valor, la funcion empalma en \(3\).
\end{fullsolution}

\section{Modelo C09.S07 --- Optimizacion}

\begin{exambrief}{Mini-examen de cierre}
Duracion sugerida: \(35\) minutos.
\begin{enumerate}
  \item Un rectangulo tiene perimetro \(24\) cm. Expresa el area en funcion de un lado.
  \item Halla las dimensiones que maximizan el area.
\end{enumerate}
\end{exambrief}

\begin{shortanswer}{Mini-examen C09.S07}
Si un lado es \(x\), el otro vale \(12-x\) y
\[
A(x)=x(12-x)=12x-x^2.
\]
El area maxima se alcanza en \(x=6\), es decir, en el cuadrado de lado \(6\).
\end{shortanswer}

\begin{fullsolution}{Mini-examen C09.S07}
Del perimetro
\[
2x+2y=24
\Longrightarrow
y=12-x.
\]
Entonces
\[
A(x)=xy=x(12-x)=12x-x^2.
\]
Derivando,
\[
A'(x)=12-2x.
\]
El punto critico cumple
\[
12-2x=0 \Longrightarrow x=6.
\]
Como
\[
A''(x)=-2<0,
\]
el punto corresponde a un maximo. Luego
\[
y=12-6=6.
\]
La figura de area maxima es un cuadrado de lado \(6\) cm.
\end{fullsolution}

\section{Modelo C10.S03 --- Regresion, residuos y prediccion responsable}

\begin{exambrief}{Mini-examen de cierre}
Duracion sugerida: \(30\) minutos.
\begin{enumerate}
  \item Un ajuste lineal es \(\hat y=1.8x+0.4\). Predice \(y\) para \(x=6\) e interpreta la
  pendiente.
  \item Si para \(x=6\) se observa \(y=12\), calcula el residuo.
  \item Explica por que un valor \(R^2=0.96\) no autoriza a extrapolar hasta \(x=100\).
\end{enumerate}
\end{exambrief}

\begin{rubricbox}{C10.S03}
\begin{itemize}
  \item Prediccion e interpretacion de la pendiente: \(4\) puntos.
  \item Residuo con signo correcto: \(3\) puntos.
  \item Juicio razonado sobre extrapolacion: \(3\) puntos.
\end{itemize}
\end{rubricbox}

\begin{shortanswer}{Mini-examen C10.S03}
La prediccion es \(11.2\), la pendiente representa \(1.8\) unidades de \(y\) por cada unidad
adicional de \(x\), y el residuo es \(12-11.2=0.8\). El valor de \(R^2\) solo describe el ajuste
en el rango observado.
\end{shortanswer}

\begin{fullsolution}{Mini-examen C10.S03}
Sustituyendo \(x=6\),
\[
\hat y=1.8\cdot6+0.4=11.2.
\]
La pendiente indica que el modelo incrementa la respuesta prevista en \(1.8\) unidades por cada
unidad adicional de \(x\). El residuo es observado menos predicho:
\[
e=y-\hat y=12-11.2=0.8.
\]
Aunque \(R^2=0.96\) indique un ajuste alto en los datos empleados, no aporta evidencia sobre el
comportamiento del proceso lejos del intervalo observado. Predecir en \(x=100\) seria una
extrapolacion y necesita una justificacion externa al valor de \(R^2\).
\end{fullsolution}

\chapter{Examenes Modelo por Bloques}

\section{Bloque B1 --- Numeros reales y algebra inicial}

\begin{exambrief}{Examen modelo}
Duracion sugerida: \(95\) minutos.
\begin{enumerate}
  \item Convierte \(0.2\overline{7}\) en fraccion generatriz.
  \item Simplifica \(\sqrt{75}-\sqrt{12}\).
  \item Halla \(a\) para que \(x-1\) sea factor de \(x^3+ax^2-4x+3\).
  \item Factoriza \(x^2-5x+6\) y resuelve \(\dfrac{x^2-5x+6}{x-2}=0\).
\end{enumerate}
\end{exambrief}

\begin{shortanswer}{B1}
\[
0.2\overline{7}=\frac{5}{18},\qquad
\sqrt{75}-\sqrt{12}=3\sqrt{3},\qquad
a=0,
\]
y la ecuacion racional equivale a \(x=3\), con \(x\neq 2\).
\end{shortanswer}

\begin{fullsolution}{B1}
Sea \(x=0.2\overline{7}\). Entonces
\[
10x=2.\overline{7},\qquad 100x=27.\overline{7}.
\]
Restando:
\[
90x=25 \Longrightarrow x=\frac{25}{90}=\frac{5}{18}.
\]
\[
\sqrt{75}-\sqrt{12}=5\sqrt{3}-2\sqrt{3}=3\sqrt{3}.
\]
Para el factor:
\[
1+a-4+3=0 \Longrightarrow a=0.
\]
Y
\[
x^2-5x+6=(x-2)(x-3).
\]
La ecuacion
\[
\frac{(x-2)(x-3)}{x-2}=0
\]
exige \(x\neq 2\) y deja como unica solucion valida \(x=3\).
\end{fullsolution}

\section{Bloque B2 --- Ecuaciones, sistemas e inecuaciones}

\begin{exambrief}{Examen modelo}
Duracion sugerida: \(100\) minutos.
\begin{enumerate}
  \item Resuelve \(\sqrt{x+5}=x-1\).
  \item Resuelve el sistema
  \[
  \begin{cases}
  x+y=7,\\
  2x-y=5.
  \end{cases}
  \]
  \item Resuelve \((x-1)(x-4)\leq 0\).
  \item Representa la solucion de
  \[
  \begin{cases}
  x\geq 0,\\
  y\geq 0,\\
  x+y\leq 5.
  \end{cases}
  \]
\end{enumerate}
\end{exambrief}

\begin{shortanswer}{B2}
\[
x=4,\qquad (x,y)=(4,3),\qquad x\in[1,4],
\]
y la region plana es el triangulo con vertices \((0,0)\), \((5,0)\) y \((0,5)\).
\end{shortanswer}

\begin{fullsolution}{B2}
La ecuacion radical exige \(x\geq 1\). Al cuadrar:
\[
x+5=(x-1)^2=x^2-2x+1
\Longrightarrow x^2-3x-4=0.
\]
Factorizando:
\[
(x-4)(x+1)=0.
\]
La unica solucion admisible es \(x=4\).

Para el sistema, sumamos:
\[
3x=12 \Longrightarrow x=4,\qquad y=3.
\]

La inecuacion cuadratica se resuelve por tabla de signos:
\[
(x-1)(x-4)\leq 0 \Longrightarrow x\in[1,4].
\]
En el plano, la tercera restriccion delimita el semiplano bajo la recta \(x+y=5\), y junto con
los ejes positivos forma el triangulo indicado.
\end{fullsolution}

\section{Bloque B3 --- Trigonometria y vectores}

\begin{exambrief}{Examen modelo}
Duracion sugerida: \(100\) minutos.
\begin{enumerate}
  \item Convierte \(\frac{5\pi}{6}\) radianes a grados.
  \item En un triangulo rectangulo con hipotenusa \(10\) y angulo agudo \(30^\circ\), halla los
  catetos.
  \item Resuelve \(\cos x=\frac12\) en \([0,2\pi)\).
  \item Dados \(\vec u=(3,4)\) y \(\vec v=(1,-2)\), calcula \(\vec u+\vec v\) y
  \(\vec u\cdot \vec v\).
\end{enumerate}
\end{exambrief}

\begin{shortanswer}{B3}
\[
\frac{5\pi}{6}=150^\circ,\qquad
\text{catetos } 5 \text{ y } 5\sqrt{3},
\]
\[
x=\frac{\pi}{3}\ \text{o}\ \frac{5\pi}{3},\qquad
\vec u+\vec v=(4,2),\quad \vec u\cdot \vec v=-5.
\]
\end{shortanswer}

\begin{fullsolution}{B3}
\[
\frac{5\pi}{6}\cdot \frac{180^\circ}{\pi}=150^\circ.
\]
En el triangulo rectangulo:
\[
\text{cateto opuesto}=10\sen 30^\circ=5,
\qquad
\text{cateto contiguo}=10\cos 30^\circ=5\sqrt{3}.
\]
La ecuacion trigonometrica se resuelve en los cuadrantes I y IV:
\[
x=\frac{\pi}{3},\qquad x=\frac{5\pi}{3}.
\]
Y en vectores:
\[
\vec u+\vec v=(3,4)+(1,-2)=(4,2),
\]
\[
\vec u\cdot \vec v=3\cdot 1+4\cdot(-2)=-5.
\]
\end{fullsolution}

\section{Bloque B4 --- Geometria analitica}

\begin{exambrief}{Examen modelo}
Duracion sugerida: \(105\) minutos.
\begin{enumerate}
  \item Halla la ecuacion de la recta que pasa por \(A(1,2)\) y \(B(5,6)\).
  \item Halla la recta perpendicular a la anterior que pasa por \(P(0,3)\).
  \item Calcula el area del triangulo de vertices \((0,0)\), \((4,0)\) y \((1,3)\).
\end{enumerate}
\end{exambrief}

\begin{shortanswer}{B4}
La recta \(AB\) es \(y=x+1\), la perpendicular por \(P\) es \(y=-x+3\) y el area del
triangulo es \(6\).
\end{shortanswer}

\begin{fullsolution}{B4}
La pendiente de \(AB\) es
\[
m=\frac{6-2}{5-1}=1.
\]
Usando el punto \(A(1,2)\),
\[
y-2=1(x-1) \Longrightarrow y=x+1.
\]
La recta perpendicular tiene pendiente \(-1\) y pasa por \(P(0,3)\):
\[
y-3=-1(x-0) \Longrightarrow y=-x+3.
\]
Para el area usamos base y altura:
\[
Area=\frac{4\cdot 3}{2}=6.
\]
\end{fullsolution}

\section{Bloque B5 --- Limites, continuidad y asintotas}

\begin{exambrief}{Examen modelo}
Duracion sugerida: \(100\) minutos.
\begin{enumerate}
  \item Halla el dominio de \(f(x)=\dfrac{\sqrt{x-1}}{x+2}\).
  \item Calcula \(\lim_{x\to 2}\dfrac{x^2-4}{x-2}\).
  \item Estudia la continuidad en \(x=1\) de
  \[
  g(x)=
  \begin{cases}
    x+1, & x<1,\\
    3x-1, & x\geq 1.
  \end{cases}
  \]
\end{enumerate}
\end{exambrief}

\begin{shortanswer}{B5}
\[
D_f=[1,\infty),\qquad
\lim_{x\to 2}\frac{x^2-4}{x-2}=4,
\]
y \(g\) es continua en \(x=1\).
\end{shortanswer}

\begin{fullsolution}{B5}
La raiz exige \(x-1\geq 0\), luego \(x\geq 1\). La restriccion \(x\neq -2\) ya queda fuera de
ese intervalo, de modo que
\[
D_f=[1,\infty).
\]
Factorizando:
\[
\frac{x^2-4}{x-2}=\frac{(x-2)(x+2)}{x-2}=x+2 \quad (x\neq 2),
\]
por lo que el limite vale \(4\).

En la funcion a trozos:
\[
\lim_{x\to 1^-}g(x)=1+1=2,
\qquad
g(1)=3\cdot 1-1=2,
\qquad
\lim_{x\to 1^+}g(x)=2.
\]
Los tres valores coinciden, asi que \(g\) es continua en \(x=1\).
\end{fullsolution}

\section{Bloque B6 --- Derivadas y aplicaciones}

\begin{exambrief}{Examen modelo}
Duracion sugerida: \(110\) minutos.
\begin{enumerate}
  \item Deriva \(f(x)=x^3-3x^2+2\).
  \item Estudia sus puntos criticos.
  \item Halla la recta tangente en \(x=1\).
\end{enumerate}
\end{exambrief}

\begin{shortanswer}{B6}
\[
f'(x)=3x^2-6x=3x(x-2).
\]
Hay un maximo local en \((0,2)\) y un minimo local en \((2,-2)\). La tangente en \(x=1\) es
\(y=-3x+3\).
\end{shortanswer}

\begin{fullsolution}{B6}
\[
f'(x)=3x^2-6x=3x(x-2).
\]
Los puntos criticos son
\[
x=0,\qquad x=2.
\]
Como \(f'\) cambia de positivo a negativo en \(x=0\), existe un maximo local en
\((0,f(0))=(0,2)\). En \(x=2\), \(f'\) cambia de negativo a positivo, por lo que hay un minimo
local en \((2,f(2))=(2,-2)\).
Ademas,
\[
f(1)=1-3+2=0,
\qquad
f'(1)=3-6=-3.
\]
La recta tangente en \(x=1\) es
\[
y-f(1)=f'(1)(x-1),
\]
es decir,
\[
y=-3(x-1)=-3x+3.
\]
\end{fullsolution}

\section{Bloque B7 --- Estadistica, probabilidad e inferencia}

\begin{exambrief}{Examen modelo}
Duracion sugerida: \(110\) minutos. Se permite calculadora o software estadistico en el segundo
apartado.
\begin{enumerate}
  \item En una clase, entre quienes usan una plataforma de estudio superan una prueba \(18\) de
  \(24\); entre quienes no la usan, \(8\) de \(16\). Calcula las dos proporciones condicionadas
  y valora si los datos prueban causalidad.
  \item Para los pares \((1,2),(2,4),(3,6),(4,8),(5,10),(6,20)\), obtiene la recta de regresion,
  \(r\) y \(R^2\). Explica el efecto del ultimo punto.
  \item Se elige al azar un equipo de \(3\) personas entre \(8\). Calcula la probabilidad de que
  dos personas concretas pertenezcan ambas al equipo.
  \item El \(4\%\) de las piezas tiene defecto. Un sistema alerta en el \(95\%\) de las
  defectuosas y en el \(3\%\) de las correctas. Calcula \(P(A)\) y \(P(D\mid A)\).
\end{enumerate}
\end{exambrief}

\begin{shortanswer}{B7}
Las proporciones son \(0.75\) y \(0.50\), pero la asociacion no prueba causalidad. La regresion
es \(\hat y\approx-2.667+3.143x\), con \(r\approx0.922\) y \(R^2\approx0.850\). La probabilidad
del equipo es \(3/28\). Finalmente, \(P(A)=0.0668\) y \(P(D\mid A)\approx0.5689\).
\end{shortanswer}

\begin{fullsolution}{B7}
Las proporciones condicionadas son
\[
P(S\mid U)=\frac{18}{24}=0.75,
\qquad
P(S\mid\overline U)=\frac8{16}=0.50.
\]
La diferencia muestra asociacion en la muestra, pero no controla motivacion, tiempo de estudio u
otras variables de confusion; por tanto, no demuestra causalidad.

Para la regresion, \(\bar x=3.5\), \(\bar y=25/3\),
\(S_{xx}=17.5\) y \(S_{xy}=55\). Por ello,
\[
b=\frac{S_{xy}}{S_{xx}}=\frac{22}{7}\approx3.143,
\qquad
a=\bar y-b\bar x=-\frac83\approx-2.667.
\]
El calculo tecnologico da \(r\approx0.922\) y \(R^2\approx0.850\). Los cinco primeros puntos
siguen exactamente \(y=2x\); el ultimo es influyente y cambia de forma notable la pendiente y el
ajuste, aunque \(r\) siga siendo alto.

Hay \(\binom83=56\) equipos posibles. Si se fijan las dos personas, quedan \(6\) opciones para
la tercera:
\[
P(\text{ambas})=\frac6{56}=\frac3{28}.
\]
Por probabilidad total y Bayes,
\[
P(A)=0.04\cdot0.95+0.96\cdot0.03=0.0668,
\]
\[
P(D\mid A)=\frac{0.04\cdot0.95}{0.0668}\approx0.5689.
\]
La alerta eleva mucho la probabilidad de defecto, pero no la convierte en certeza.
\end{fullsolution}

\chapter{Hoja de Ruta de Produccion}

\begin{planningbox}{Orden recomendado}
\begin{enumerate}
  \item Redactar primero todos los mini-examenes por seccion usando una plantilla fija.
  \item Construir despues los examenes por bloques a partir de combinaciones ya probadas.
  \item Anadir baremos detallados en una segunda pasada, cuando el peso de cada item este
  equilibrado.
  \item Cerrar con simulacros largos y control de tiempos reales.
\end{enumerate}
\end{planningbox}

\begin{databox}{Criterios de calidad del suplemento definitivo}
\begin{itemize}
  \item cada mini-examen debe poder resolverse sin mirar teoria, pero exigiendo reconocer el
  apartado;
  \item cada examen por bloques debe incluir al menos un item de decision metodologica;
  \item ningun bloque debe quedarse sin solucion completa ni sin baremo;
  \item las figuras y tablas deben ser legibles en impresion y coherentes en escala de grises.
\end{itemize}
\end{databox}
